%La informaci\'on que describen en este punto se basa en los estudios realizados con antelaci\'on respecto al tema.
%Lo m\'as importante que se ha hecho hasta ahora con relaci\'on al problema. Deber\'ia responder preguntas como las siguientes:
%?`cu\'ando surge?, ?`qu\'e m\'etodos se han usado para resolverlo?, ?`cu\'ales son los mejores algoritmos que se han creado hasta
%la fecha?, ?`qu\'e representaciones han tenido los mejores resultados?, ?`cu\'al es la tendencia actual para resolver el problema?, tipos de movimientos, heur\'isticas, m\'etodos completos, tendencias, etc... Puede incluir gr\'aficos comparativos o explicativos.\\

\subsection{Surgimiento}

Los drones han sido utilizados por años en trabajos militares, pero no fue hasta el 2006 cuando se permitió por primera vez en Estados Unidos utilizar drones para fines comerciales. Tras esto, ha habido un auge en la utilización de los drones. En \textit{Moshref-Javadi} y \textit{Winkenbach} \cite{moshrefjavadi2021applications}, se tiene que ya se utilizan drones para servicios de comercio (como \textit{Amazon} o \textit{Walmart}), en servicios de alimentación (como \textit{Coca-Cola} o \textit{Domino's}), e incluso han sido utilizados en servicios de emergencia. Por ejemplo, en \textit{Stell et al.} \cite{stell2022drone} se menciona la utilización de drones para hacer entregas de medicamentos durante el periodo del COVID-19.

Según \cite{transinfo_drones}, los drones han surgido de manera explosiva en la logística debido a su capacidad de bajar los costos de entrega. Efectivamente, estudios como \cite{moshrefjavadi2021applications} explican que el costo de entrega por paquete puede disminuir a $0.88$ dólares por entrega, a diferencia de utilizar únicamente camiones con un costo aproximado de $5.99$ dólares.

El trabajo seminal de \textit{Murray} y \textit{Chu} \cite{Murray2015Flying} es considerado una de las primeras referencias con respecto al problema en un formato de optimización. Hay varios algoritmos que utilizan este trabajo como una base para sus propias propuestas, tales como los trabajos de \textit{Ponza} \cite{ponza2016optimization} y \textit{Freitas} y \textit{Penna} \cite{freitas2021exact}.

\subsection{Métodos de Resolución}

El trabajo de \textit{Murray} y \textit{Chu} ya hace referencia a la complejidad del problema. Como se mencionó previamente, el problema se considera \textit{NP-hard}, por lo que no se puede resolver de forma exacta en tiempo polinomial. Frente a esto, se han utilizado varias estrategias según el trabajo considerado.

Se han realizado enfoques de técnicas de resolución completas. En \textit{Murray} y \textit{Chu} \cite{Murray2015Flying}, se utilizan solvers como \textit{CPLEX} y \textit{Gurobi}, los cuales terminan de saturar el problema con una cantidad de $10$ clientes. En \textit{Freitas} y \textit{Penna} \cite{freitas2021exact}, se utiliza una indexación de los momentos del camión en un conjunto $L$, lo cual permite eliminar el tiempo de computación que se utiliza en el trabajo de \textit{Murray} y \textit{Chu}, dado que este requería eliminar subtours mediante una variable auxiliar $u_i$. En \textit{Moshref-Javadi} y \textit{Winkenbach} \cite{moshrefjavadi2021applications} se menciona la utilización de algoritmos de \textit{Branch-and-Cut}.

% Notas del ayudante
%Se han realizado enfoques de técnicas metaheurísticas. En \cite{freitas2021exact} se hace uso de una heurística constructiva haciendo una descomposición del problema asignando inicialmente las entregas dentro del radio del dron para dejar las demás al camión. Se hace referencia a enfoques de resolución genéticos en \cite{moshrefjavadi2021applications}.

Hay trabajos que han realizado enfoques de técnicas metaheurísticas. En \textit{Freitas} y \textit{Penna} \cite{freitas2021exact} se hace con construcción de ruta inicial. Esta se implementa mediante un algoritmo de \textit{Tabu Search} mejorado, definido como \textit{Hybrid Tabu General Variable Neighborhood Search}. Se parte de un tour terrestre óptimo resuelto mediante \textit{Concorde TSP} y se genera una solución inicial mediante ahorros netos. Su método combina la exploración de vecindarios variables (\textit{RVND}) con una lista de memoria tabú que prohíbe que clientes transferidos al dron retornen inmediatamente al camión. Según su informe, superó las mejores soluciones conocidas en más del $80\%$ de los casos de prueba.

%En \cite{ponza2016optimization} se implementó una metaheurística basada en el algoritmo de Monte Carlo (haciendo intercambios entre rutas terrestres y aéreas). En \cite{freitas2021exact} se desarrolló un algoritmo propio basado en la técnica de memoria Tabu, construyendo una ruta inicial y trabajando sobre esta, haciendo intercambios entre entregas de dron y entregas de camión. Este utiliza la lista Tabu para evitar los ciclos.

En \textit{Ponza} \cite{ponza2016optimization} se implementó \textit{Simulated Annealing}. Este trabajo es el primero en indicar pruebas entre $5$ y $200$ clientes, el cual es un estándar para las pruebas que se define en la mayoría de los artículos. En el presente trabajo también se hace referencia a este rango. Utilizando una representación eficiente en arreglos contiguos y operadores de reubicación de clientes (\textit{relocate}), se evaluó el intercambio dinámico entre rutas terrestres y salidas aéreas $\langle i, j, k \rangle$, descartando movimientos infactibles respecto a la autonomía de batería $E$.

%Además, se han hecho transformaciones del problema para trabajarlo de manera equivalente. En \cite{boysen2018drone}, se propone fijar la ruta del camión y utilizar un único dron con la condición de volver a la misma parada de lanzamiento. Esta forma del problema se puede resolver en tiempo polinomial usando un algoritmo voraz. Agregar que el dron debe volver en la misma parada o en una consecutiva también puede resolverse en tiempo polinomial haciendo uso de un algoritmo de asignación lineal.

Sobre estas optimizaciones, se han hecho análisis para determinar el tiempo de cómputo bajo ciertas condiciones del camión. En particular, \textit{Boysen et al.} \cite{boysen2018drone} formalizaron el \textit{Drone Scheduling Problem} (DSP) para una ruta fija de camión, caracterizando su complejidad teórica según la política de aterrizaje: si el dron debe retornar a la misma parada de lanzamiento (Política A), el subproblema se resuelve óptimamente en tiempo polinomial $\mathcal{O}(nK)$ mediante una asignación voraz; si se le permite retornar en la misma parada o en la inmediatamente consecutiva (Política B), demostraron que se modela y resuelve en tiempo polinomial como un \textit{Linear Assignment Problem} (LAP).

% Sección extra del evolutivo.

Por último, con énfasis en la implementación genética y evolutiva, se hacen referencias a estas técnicas en varios trabajos. En \textit{Ha et al.} \cite{ha2020hybrid} se realiza una implementación genética híbrida. El cromosoma se representa como un tour TSP y se utiliza un operador de cruce especializado para el dron, enfocándose en subrutas aéreas. En \textit{Pambo et al.} \cite{pambo2024evolutionary} se hace uso de separar la información de los nodos visitados y su orden, junto a un arreglo indicando si el cliente fue atendido por camión o dron. Se implementa cruzamiento, mutación y elitismo, haciendo mayor énfasis en cruzamiento. Asimismo, incorpora una función reparadora en caso de que se asignen vuelos infactibles.

\subsection{Tendencias}

La tendencia se orienta a resolver el problema mediante heurísticas incompletas. En \textit{Moshref-Javadi} y \textit{Winkenbach} \cite{moshrefjavadi2021applications} se hace una comparación de las técnicas utilizadas por diversos artículos científicos, concluyendo que un $53\%$ recurre a heurísticas y metaheurísticas para la resolución del problema. En general, los algoritmos completos se reservan para instancias pequeñas (no más de $10$ clientes), mientras que las heurísticas y metaheurísticas se utilizan con valores cercanos a $n=200$.

% Se me olvido guardar el original. Modificado de aca en adelante

Del trabajo de \textit{Moshref-Javadi} y \textit{Winkenbach} \cite{moshrefjavadi2021applications} se concluye que la técnica principal se conoce como el enfoque \textit{Truck-First, Drone-Second}, el cual se utiliza en el presente algoritmo evolutivo. Consiste en determinar primero el recorrido del camión (es decir, resolver un problema de ruteo base) y luego ir realizando modificaciones dentro de este camino para agregarle entregas de dron. La idea es identificar qué trayectos del dron permiten disminuir el tiempo total mediante algoritmos constructivos, teniendo la solución inicial y mejorándola de forma iterativa.

Implementar de esta manera el evolutivo tiene sus ventajas según la literatura. En \textit{Prins} \cite{prins2004simple} se explica cómo la factibilidad siempre se garantiza, debido a que todo recorrido mixto puede realizarse enteramente por el camión. Además, está la ventaja de la velocidad de cómputo mencionada en \textit{Boysen et al.} \cite{boysen2018drone}. El principal desafío señalado en \textit{Boysen et al.} \cite{boysen2018drone} radica en que una ruta óptima para el TSP no necesariamente estará optimizada para el dron, lo que induce cierta miopía en la solución.
